\documentclass[a4paper,12pt]{article}
\usepackage[utf8]{inputenc}
\usepackage[T1]{fontenc}
\usepackage[english]{babel}
\usepackage{geometry}
\geometry{left=2.5cm,right=2.5cm,top=2.5cm,bottom=2.5cm}

\usepackage{lmodern}
\usepackage{xcolor}
\definecolor{titleblue}{RGB}{0,51,140}

\usepackage{hyperref}

\pagestyle{empty}

\begin{document}

\begin{flushleft}
    \textbf{Ismail AISSAOUI} \\
    State Engineer in Artificial Intelligence \\
    ismail.aissaoui.pro@gmail.com \\
    +213 660 70 77 96 \\
    \href{https://ismail-aissaoui.vercel.app}{ismail-aissaoui.vercel.app}
\end{flushleft}

\begin{flushright}
    To the PhD Selection Committee \\
    \textbf{EDT Program - PC2} \\
    Attn: Johann Bourcier, Benoît Combemale \\
    \medskip
    May 3, 2026
\end{flushright}

\bigskip

\noindent \textbf{Subject: Application for PhD: Composing and Configuring Digital Twin Architectures to Meet User Requirements (Rank 9 - PC2)}

\bigskip

Dear Members of the Selection Committee,

As a State Engineer in Artificial Intelligence from ESI-SBA, I am writing to express my strong interest in the PhD position "Composing and configuring digital twin architectures to meet user requirements (energy, performance, budget)" within the Engineering Digital Twin (EDT) national program. This application is motivated by the alignment between my background in system optimization and the research objectives of your group.

The challenge of navigating a **Digital Twin Architectural Design Space (DT-ADS)** to balance competing non-functional requirements is at the heart of my current research. In my Master's thesis at Sonatrach, I designed a **Physics-Informed Digital Twin** for gas turbines where I encountered the "Latency Paradox." To resolve this, I implemented a hardware-software co-design strategy, optimizing a hybrid **Mamba-SSM** and **xLSTM** architecture to achieve **0.04ms inference** on edge devices. This experience directly addresses your objective of linking architectural choices to constraints like performance and energy consumption.

Furthermore, I have developed extensive expertise in **automated design space exploration**. In my "Advanced NLP" project, I built a platform to optimize DistilBERT models using nature-inspired metaheuristics, including **Particle Swarm Optimization (PSO)**, **Genetic Algorithms (GA)**, and **Bayesian Optimization**. I am eager to apply these optimization paradigms to the modeling of DT-ADS, facilitating the automated composition of complex Digital Twin pipelines—from sensors to decision modules.

My technical stack—ranging from high-performance **C++ JIT** and **PyTorch** to distributed synchronization logic in **PostgreSQL/SQLite** systems—positions me to contribute significantly to the **Artemis platform**. I am a highly autonomous researcher with a deep commitment to open-source software and industrial innovation.

Thank you for your time and consideration. I look forward to discussing how my experience in architectural optimization and hybrid modeling can contribute to the success of the PC2 axis and the broader EDT ecosystem.

Sincerely,

\bigskip

\textbf{Ismail AISSAOUI}

\end{document}
